\section{Scope}\label{a:sec:introduction}

For the subcritical binary Galton--Watson process, late survival has the form
\[
 S_n\sim A(p)(2p)^n,
 \qquad 0\le p<\tfrac12,
\]
and the limiting population mean conditional on survival is \(1/A(p)\).  The
probabilistic construction, extinction theory and conditional-moment calculation
were developed in \ChM.  The present chapter takes the resulting constant as its
object and develops its analytic structure.

The argument is developed in layers.  An exact infinite product leads to a
telescoping series, two-sided and parity bounds, and a near-critical theorem with a
controlled remainder.  These results make it possible to compare the discrete
constant with the explicit continuous-time value \(A_{\mathrm c}(p)\) before
assessing the negative closed-form search.  The second half then identifies
\(A(p)\) with a basin value of the Koenigs linearising function for the logistic
map.  The Becker--Bergweiler classification justifies a precise obstruction: for
each fixed \(r\in(0,1)\), the function \(z\mapsto\psi_r(z)\) is
hypertranscendental.  It does not determine the arithmetic nature of any particular
number \(A(p_0)\), or the differential algebraicity of the parameter map
\(p\mapsto A(p)\).  The scope taxonomy and inverse-symbolic evidence are organised
to keep those questions distinct.

The chapter closes by turning the analytic results into a practical evaluation
scheme.  The elementary repelling cases \(r=2\) and \(r=4\), both outside the
subcritical window, are derived in the appendix.
